\documentclass[12pt, a4paper]{article}

% ─── Packages ─────────────────────────────────────────────────────────────────
\usepackage[utf8]{inputenc}
\usepackage[T1]{fontenc}
\usepackage[french]{babel}
\usepackage{amsmath, amssymb}
\usepackage{geometry}
\usepackage{hyperref}
\usepackage{booktabs}
\usepackage{array}
\usepackage{longtable}
\usepackage{graphicx}
\usepackage{cite}
\usepackage{setspace}
\usepackage{parskip}
\usepackage{titlesec}
\usepackage{abstract}
\usepackage{fancyhdr}
\usepackage{xcolor}

\geometry{top=2.5cm, bottom=2.5cm, left=2.5cm, right=2.5cm}
\onehalfspacing

\hypersetup{
    colorlinks=true,
    linkcolor=blue!60!black,
    citecolor=blue!60!black,
    urlcolor=blue!60!black
}

% ─── En-tête ─────────────────────────────────────────────────────────────────
\pagestyle{fancy}
\fancyhf{}
\rhead{\small Détection de variations atypiques par IA  Cameroun / BEAC}
\lhead{\small État de l'art}
\rfoot{\thepage}

% ─── Titre ───────────────────────────────────────────────────────────────────
\title{
    \textbf{Détection de Variations Atypiques dans les Séries Temporelles\\
    Financières par l'Intelligence Artificielle :}\\[0.3em]
    \large État de l'Art
}

\author{
    \textbf{MOUKOKO EKAMBI GEORGES MIGUEL}\\[0.3em]
    \textit{Stagiaire, Banque des États de l'Afrique Centrale}\\
    \textit{M2 Data Science \& IA  ENSPD / Université de Douala}\\
    \textit{LIDIA Laboratory}\\[0.5em]
    \textit{Données sources : Banque des États de l'Afrique Centrale (BEAC)}\\
}

\date{\today}

% ─── Document ────────────────────────────────────────────────────────────────
\begin{document}

\maketitle

% ─── Résumé ──────────────────────────────────────────────────────────────────
\begin{abstract}
La détection de variations atypiques  aussi appelée \textit{anomaly detection}
 constitue un enjeu central de la surveillance des systèmes financiers.
Ce document dresse un état de l'art des principales méthodes d'intelligence
artificielle mobilisées à cette fin, depuis les approches fondatrices basées
sur la densité locale jusqu'aux architectures profondes les plus récentes
(autoencodeurs, réseaux génératifs adversariaux, transformers).
Chaque méthode est présentée avec ses fondements théoriques, ses forces et ses
limites, en s'appuyant sur les travaux scientifiques de référence.
Nous présentons également \textbf{MKO-KAN}, une architecture originale à base
de réseaux de Kolmogorov-Arnold développée dans le cadre de ce stage, qui apporte
une capacité d'interprétabilité symbolique adaptée aux exigences d'audit COBAC/BEAC.
En seconde partie, nous présentons le jeu de données Monétaires
produit à partir des statistiques de la BEAC (Banque des États de
l'Afrique Centrale) selon la nomenclature IFS du Fonds Monétaire International,
et nous analysons la cohérence entre les caractéristiques de ce dataset et
chacune des méthodes IA exposées, afin d'orienter les choix méthodologiques
pour la détection d'anomalies dans ce contexte.
\end{abstract}

\tableofcontents
\newpage

% ═════════════════════════════════════════════════════════════════════════════
\section{Introduction}
% ═════════════════════════════════════════════════════════════════════════════

La détection de variations atypiques est définie comme le processus
d'identification de patrons dans des données qui ne se conforment pas au
comportement attendu. Dans le domaine financier et monétaire, ces variations
peuvent signaler des erreurs de collecte, des chocs macroéconomiques, des
fraudes, ou des transitions structurelles importantes.
L'essor des méthodes d'apprentissage automatique a profondément renouvelé
l'approche de ce problème, en permettant de traiter des données de haute
dimension, des séries temporelles longues et des structures multivariées
complexes.

Comme le soulignent Pang \textit{et al.} \cite{pang2021deep}, la détection
d'anomalies par apprentissage profond constitue désormais une direction
critique, couvrant trois grandes catégories de méthodes et onze sous-catégories
de techniques, chacune reposant sur des hypothèses et des fonctions objectif
spécifiques.

Cette étude est structuré comme suit : la section~\ref{sec:methodes_classiques}
présente les méthodes classiques de référence ; la section~\ref{sec:deep}
couvre les approches d'apprentissage profond, dont MKO-KAN ; la
section~\ref{sec:dataset} décrit le jeu de données BEAC-Cameroun ; enfin,
la section~\ref{sec:coherence} évalue la cohérence méthodologique entre chaque
technique et les caractéristiques du dataset.


% ═════════════════════════════════════════════════════════════════════════════
\section{Méthodes Classiques de Référence}
\label{sec:methodes_classiques}
% ═════════════════════════════════════════════════════════════════════════════

\subsection{Local Outlier Factor (LOF)}

Le \textit{Local Outlier Factor} (LOF) est l'une des méthodes fondatrices
de la détection d'anomalies fondées sur la densité locale, introduite par
Breunig \textit{et al.} \cite{breunig2000lof}. La contribution centrale de
ce travail réside dans l'attribution à chaque observation d'un \textit{degré}
d'anormalité  le facteur LOF  plutôt qu'une simple classification
binaire. Ce degré est calculé en comparant la densité locale d'un point à
celle de ses $k$ plus proches voisins.

Formellement, la densité locale de chaque observation $p$ est estimée par la
distance de joignabilité (\textit{reachability distance}) à ses $k$ voisins.
Un point dont la densité locale est significativement inférieure à celle de
son voisinage se voit attribuer un LOF élevé, le qualifiant d'anomalie.

Cette approche présente l'avantage d'être non-paramétrique et de s'adapter
aux distributions de données complexes et multi-modales. Ses limites tiennent
à son coût computationnel en $\mathcal{O}(n^2)$ et à sa sensibilité au
paramètre $k$.




\subsection{Isolation Forest (iForest)}

L'Isolation Forest, proposé par Liu \textit{et al.} \cite{liu2012isolation},
constitue une rupture paradigmatique par rapport aux méthodes précédentes.
Plutôt que de modéliser le comportement normal pour en déduire les anomalies,
iForest isole directement les anomalies en exploitant deux propriétés
intrinsèques : elles sont en minorité dans le dataset et leurs attributs
diffèrent significativement des observations normales.

L'algorithme construit un ensemble (\textit{forest}) d'arbres d'isolation
(\textit{isolation trees}), chaque arbre partitionnant récursivement l'espace
des données par des coupures aléatoires. La longueur moyenne du chemin
parcouru dans l'arbre pour isoler un point est inversement proportionnelle
à son score d'anomalie : les points anormaux sont isolés plus rapidement car
ils résident dans des régions de faible densité.

Liu \textit{et al.} \cite{liu2012isolation} démontrent qu'iForest surpasse les
méthodes One-Class SVM, LOF et ORCA en termes d'AUC et de temps de traitement,
avec une complexité linéaire $\mathcal{O}(n)$ particulièrement adaptée aux
grandes volumétries. Il s'avère également robuste aux phénomènes de masquage
et de dilution, et performant en haute dimension.



% ═════════════════════════════════════════════════════════════════════════════
\section{Méthodes d'Apprentissage Profond}
\label{sec:deep}
% ═════════════════════════════════════════════════════════════════════════════

\subsection{Autoencodeurs (Autoencoder, AE)}

Les autoencodeurs constituent la porte d'entrée du \textit{deep learning}
dans la détection d'anomalies. Introduits pour la réduction de dimensionnalité
par Hinton \& Salakhutdinov \cite{hinton2006reducing}, ils se composent
d'un encodeur qui projette les données dans un espace latent de dimension
réduite et d'un décodeur qui reconstruit l'entrée originale.

Le principe central pour la détection d'anomalies repose sur l'hypothèse
de reconstruction différentielle : un modèle entraîné exclusivement sur
des données normales développe une capacité de reconstruction élevée pour
les exemples normaux, mais produit une erreur de reconstruction
(\textit{reconstruction error}) élevée pour les anomalies, dont les
patterns ne correspondent pas à la distribution apprise.

Le seuil de détection est généralement fixé sur la distribution empirique
des erreurs de reconstruction observées sur l'ensemble d'entraînement.
Parmi les variantes notables, les Robust Deep Autoencoders (RDA) décomposent
l'entrée $X$ en une composante normale $L_D$ et une composante de bruit $S$
\cite{zhou2017anomaly}, renforçant la robustesse aux contaminations.



\subsection{Autoencodeurs Variationnels (Variational Autoencoder, VAE)}

Le VAE, formalisé par Kingma \& Welling \cite{kingma2013auto}, étend
l'autoencodeur classique en introduisant une représentation probabiliste
de l'espace latent. Plutôt que d'encoder chaque entrée en un point fixe,
le VAE encode une distribution de probabilité (typiquement gaussienne)
dans l'espace latent.

Pour la détection d'anomalies, au lieu de mesurer une simple erreur de
reconstruction, on calcule une \textit{probabilité de reconstruction}
estimée par échantillonnage multiple depuis l'espace latent \cite{an2015vae}.
Cette approche probabiliste permet de mieux quantifier l'incertitude et
d'attribuer des scores d'anomalie plus nuancés.

Le VAE est particulièrement adapté à la détection d'anomalies non supervisée
dans les séries temporelles multivariées, bien que sa capacité à modéliser
les dépendances à long terme reste une limitation reconnue \cite{shap2025}.
Des architectures hybrides intégrant des mécanismes d'attention
multi-têtes ont été proposées pour pallier cette limite.



\subsection{Réseaux Génératifs Adversariaux (GAN)  AnoGAN et f-AnoGAN}

Les GANs ont ouvert une nouvelle voie pour la détection d'anomalies non
supervisée. Schlegl \textit{et al.} \cite{schlegl2017unsupervised} ont
introduit \textbf{AnoGAN}, le premier framework GAN dédié à la détection
d'anomalies, en exploitant la capacité des GANs à modéliser la distribution
des données normales.

L'architecture repose sur l'entraînement d'un GAN (architecture DCGAN)
uniquement sur des données normales, apprenant ainsi la distribution de
la variabilité normale dans l'espace latent. Pour détecter une anomalie,
AnoGAN cherche le point $z$ dans l'espace latent tel que le vecteur généré
$G(z)$ soit le plus proche de l'observation $x$ à tester. Un score d'anomalie
combinant l'erreur de reconstruction et un score de discrimination est alors
calculé. L'anomalie correspond à une observation qui ne peut pas être
fidèlement mappée dans l'espace latent du GAN.

La version améliorée \textbf{f-AnoGAN} \cite{schlegl2019fanogan} introduit
un encodeur dédié pour accélérer drastiquement l'inférence, en substituant
à la recherche itérative dans l'espace latent un mapping direct image-vers-
latent. Cette amélioration rend l'approche praticable dans des contextes
opérationnels.

\subsection{Transformers  Anomaly Transformer}

Les Transformers, initialement conçus pour le traitement du langage naturel,
ont récemment démontré un potentiel remarquable pour la détection d'anomalies
dans les séries temporelles, en exploitant les mécanismes d'auto-attention
pour capturer les dépendances à longue portée \cite{mamba2025}.

L'\textbf{Anomaly Transformer}, proposé par Xu \textit{et al.}
\cite{xu2022anomaly} (ICLR 2022, Spotlight), introduit le concept
d'\textit{Association Discrepancy} comme critère discriminant. L'observation
centrale est la suivante : en raison de la rareté des anomalies, un point
anormal ne peut construire que des associations triviales avec le reste de la
série temporelle, concentrées sur ses points temporels adjacents. Un point
normal, en revanche, construit des associations distribuées sur l'ensemble
de la série.

L'Anomaly Transformer formalise cette intuition en apprenant deux
distributions d'associations  une distribution gaussienne prioritaire
(\textit{prior association}) et une association à partir des données
(\textit{series association})  et en maximisant leur discordance via
une stratégie \textit{minimax}. Cette approche surpasse les méthodes
antérieures sur plusieurs benchmarks de séries temporelles.

Des extensions récentes intègrent des mécanismes d'attention parcimonieuse
(\textit{sparse attention}) pour améliorer l'efficacité sur les longues
séquences, ainsi que des architectures hybrides Mamba-Transformer
\cite{mamba2025} qui combinent les avantages des modèles d'état sélectifs
et de l'auto-attention.


\subsection{MKO-KAN : Réseaux de Kolmogorov-Arnold Hybrides pour la Détection
d'Anomalies Financières}

Les réseaux de Kolmogorov-Arnold (KAN), formalisés par Liu \textit{et al.}
\cite{liu2024kan}, constituent une rupture architecturale par rapport aux
perceptrons multicouches classiques (MLP). Là où un MLP applique des fonctions
d'activation fixes aux nœuds, un KAN paramétrise des \textit{fonctions d'arête
apprenables} $\varphi_{ij}$ entre chaque paire de nœuds connectés, en
s'appuyant sur le théorème de représentation de Kolmogorov-Arnold. Cette
propriété confère aux KAN une interprétabilité structurelle native : après
élagage, le réseau se réduit à un graphe de formules analytiques.

Le modèle \textbf{MKO-KAN} (\textit{Mobile-Money Kolmogorov-Arnold Network}),
développé dans le cadre de ce stage à la BEAC \cite{moukoko2025mkan}, adapte
cette architecture à la détection de fraude et d'anomalies dans les séries
financières de la zone CEMAC. Il repose sur trois innovations clés :

\begin{enumerate}
    \item \textbf{Cellule T-KAN (TKANCell)} : extension récurrente de type LSTM
    dans laquelle chaque porte (oubli, entrée, candidate, sortie) est une couche
    KAN hybride \texttt{HybridKANLayer}. Cette structure capture les dépendances
    temporelles séquentielles propres aux flux transactionnels ou aux agrégats
    monétaires mensuels.

    \item \textbf{Arêtes hybrides Gaussienne + Fourier (HybridEdgeFunction)} :
    chaque fonction d'arête combine une composante gaussienne (FastKAN
    \cite{li2024fastkan}) pour la localisation et une composante de Fourier
    (KAN-AD \cite{zhou2024kanad}) pour les patterns périodiques :
    \[
      \varphi_{ij}(x) =
        \underbrace{\sum_{m} w_m \exp\!\left(-\frac{(x-\mu_m)^2}{2h^2}\right)}_{\text{Gaussienne (FastKAN)}}
        +
        \underbrace{\sum_{k} \left[a_k \cos(kx) + b_k \sin(kx)\right]}_{\text{Fourier (KAN-AD)}}
    \]

    \item \textbf{Nœuds multiplicatifs MultKAN} : en complément des nœuds
    additifs ($\sum_i \varphi_{ij}(x_i)$), des nœuds multiplicatifs
    ($\prod_i \varphi_{ij}(x_i)$) encodent les interactions contextuelles
    ciblées entre variables  par exemple, la co-occurrence d'un retrait
    important et d'un horaire nocturne dans les séries transactionnelles
    \cite{yanglu2024kan}.
\end{enumerate}

\textbf{Interprétabilité et conformité réglementaire.}
La sparsification par norme $L_1$ élague le réseau jusqu'à un sous-graphe
minimal, sur lequel une régression symbolique extrait des formules analytiques
explicites. Ce mécanisme produit des règles de décision auditables,
directement exploitables dans le cadre des obligations COBAC/BEAC de
traçabilité algorithmique.

La fonction de perte totale combine trois termes :
\[
  \mathcal{L} = \mathcal{L}_{\text{BCE}} + \lambda \|\Phi\|_1
              + \mu_1 H_1(\Phi) + \mu_2 H_2(\Phi)
\]
où $\mathcal{L}_{\text{BCE}}$ est la perte de classification pondérée par la
prévalence de fraude, $\|\Phi\|_1$ la régularisation $L_1$ sur les arêtes,
et $H_1, H_2$ des termes d'entropie favorisant la parcimonie structurelle.

La recherche des hyperparamètres optimaux (taille de couche cachée, bases
Gaussienne/Fourier, coefficients de régularisation, taux d'apprentissage)
est assurée par un algorithme génétique avec mutation adaptative et modèle
d'estimation de distribution (EDA), inspiré des méta-heuristiques de Hao
\& Solnon \cite{haosol2024meta}.


% ═════════════════════════════════════════════════════════════════════════════
\section{Présentation du Dataset : Données Monétaires (BEAC)}
\label{sec:dataset}
% ═════════════════════════════════════════════════════════════════════════════

\subsection{Origine et nature des données}

Les jeux de données analysées dans ce travail sont issus de la
\textbf{Banque des États de l'Afrique Centrale (BEAC)}, constitué à partir
de la nomenclature IFS (\textit{International Financial Statistics}) du
Fonds Monétaire International. Il porte sur la situation de l'\textbf{actif/passif
du bilan des Autres Sociétés de Dépôts (2SR)} de la sous région, couvrant
l'ensemble des institutions de dépôts non bancaires centrales.

Le fichier source  \texttt{clean\_<pays>\_beac\_mapping\_2SR\_<actif/passif>.xlsx}
 résulte d'un travail de nettoyage et de restructuration des données
brutes de la BEAC réalisé par l'équipe de recherche.

\subsection{Caractéristiques structurelles}

Le dataset présente les caractéristiques suivantes :

\begin{longtable}{>{\bfseries}p{4.5cm} p{9cm}}
\toprule
\textbf{Dimension} & \textbf{Description} \\
\midrule
\endhead
Période couverte & Décembre 2001 à Mars 2026 (données mensuelles)\\
Nombre de périodes & $\approx 195$ observations temporelles\\
Nombre de séries & $\approx 55$ indicateurs financiers\\
Format & Données en panel large (\textit{wide format})\\
Unité & Milliards de FCFA (Franc CFA de la CEMAC)\\
Fréquence & Mensuelle (notation : \texttt{AAAMyy})\\
\bottomrule
\end{longtable}
% ═════════════════════════════════════════════════════════════════════════════
\section{Cohérence entre le Dataset et les Méthodes IA}
\label{sec:coherence}
% ═════════════════════════════════════════════════════════════════════════════

\subsection{Cohérence avec le Local Outlier Factor (LOF)}

\textbf{Applicabilité : Modérée.}

Le LOF est adapté à la détection d'anomalies ponctuelles dans un espace
multidimensionnel statique. Son application directe aux séries temporelles
nécessite une transformation préalable (ex. : fenêtres glissantes ou
extraction de caractéristiques statiques par série).

Dans le contexte BEAC, le LOF peut être appliqué sur les niveaux mensuels
de chaque indicateur, en traitant chaque observation $(t, \mathbf{x}_t)$
comme un point dans un espace de dimension égale au nombre de séries.
La nature multivariée du dataset (55 séries) rend cependant le LOF
sensible à la malédiction de la dimensionnalité (\textit{curse of
dimensionality}), nécessitant une réduction préalable de dimension (PCA,
UMAP). Par ailleurs, les valeurs manquantes doivent être imputées et
les tendances longues supprimées pour éviter des faux positifs.

\textbf{Usage recommandé} : exploration préliminaire, détection d'anomalies
sur séries individuelles après stationnarisation.

\subsection{Cohérence avec l'Isolation Forest}

\textbf{Applicabilité : Élevée.}

L'Isolation Forest constitue l'une des méthodes les plus directement
applicables au dataset BEAC pour plusieurs raisons. Sa complexité
linéaire $\mathcal{O}(n)$ le rend efficace sur les $\approx 195 \times 55$
cellules disponibles. Sa robustesse à la haute dimensionnalité (55 séries)
et aux attributs non pertinents est particulièrement précieuse dans ce
contexte où de nombreuses séries sont des composantes d'agrégats.

La nature non supervisée d'iForest s'avère également adaptée à l'absence
de labels d'anomalies historiquement certifiés. Les valeurs manquantes
devront être traitées préalablement (imputation ou masquage).

La capacité d'iForest à fonctionner sans hypothèse de distribution
lui permet de gérer les distributions asymétriques et les queues lourdes
fréquentes dans les données financières. Il sera appliqué de manière
optimale en mode multivarié sur les résidus après décomposition des
tendances.

\subsection{Cohérence avec l'Autoencoder (AE)}

\textbf{Applicabilité : Élevée, sous conditions.}

L'autoencoder est particulièrement adapté à la structure multivariée
et temporelle du dataset BEAC. Un autoencoder récurrent (LSTM-AE ou
GRU-AE) peut modéliser simultanément les dépendances temporelles et
les co-variations entre les 55 séries. Le principe de reconstruction
différentielle est cohérent avec la logique de surveillance financière :
une série qui s'écarte de son comportement habituel génère une erreur
de reconstruction élevée.

La quantité de données ($\approx 195$ points mensuels) constitue
cependant une contrainte : un AE profond nécessitera une architecture
modérée (2--3 couches) et des techniques de régularisation (dropout,
early stopping) pour éviter le surapprentissage. La normalisation des
séries (par unité de variance ou min-max) est impérative compte tenu
de l'hétérogénéité des échelles.

\textbf{Nous recommandons notamment ces variante} : Convolutional-Recurrent AE (CR-AE) pour
exploiter à la fois la structure temporelle locale et les dépendances
longues.

\subsection{Cohérence avec le VAE}

\textbf{Applicabilité : Élevée, avec avantage probabiliste.}

Le VAE présente un avantage décisif par rapport à l'AE classique dans
le contexte BEAC : sa nature probabiliste permet de quantifier
l'incertitude associée à chaque score d'anomalie, ce qui est
particulièrement précieux pour la surveillance macroprudentielle.

La capacité du VAE à modéliser une distribution sur l'espace latent
est adaptée à la nature stochastique des agrégats monétaires, soumis
à des chocs exogènes (politiques, crises). Le VAE peut de plus être
conditionné (\textit{conditional VAE}) sur des variables temporelles
(mois, trimestre, phase du cycle économique) pour mieux capturer la
saisonnalité.

La limitation en matière de dépendances longues peut être partiellement
compensée par l'intégration de mécanismes d'attention, comme proposé
dans SHAPAttenVAE \cite{shap2025}.

\subsection{Cohérence avec les GAN (AnoGAN / f-AnoGAN)}

\textbf{Applicabilité : Modérée, contexte exigeant.}

Les approches GAN présentent un potentiel élevé pour modéliser la
distribution des régimes normaux du système bancaire camerounais.
Cependant, leur applicabilité directe est contrainte par la taille
modeste du dataset ($\approx 195$ observations) : l'entraînement d'un
GAN stable nécessite généralement plusieurs milliers d'exemples.

Une stratégie viable consiste à utiliser le \textbf{TAnoGAN}
\cite{tanogan2020}, qui adapte le framework AnoGAN aux séries temporelles
avec un faible nombre d'observations, ou à recourir à des techniques
d'augmentation de données (bootstrapping saisonnier, simulation DSGE
calibrée). f-AnoGAN reste la variante la plus opérationnelle grâce à
son encodeur dédié qui élimine la phase de recherche dans l'espace latent.

\subsection{Cohérence avec l'Anomaly Transformer}

\textbf{Applicabilité : Très élevée pour les séries longues.}

L'Anomaly Transformer est particulièrement adapté à la nature temporelle
et multivariée du dataset BEAC. Son mécanisme d'\textit{Association
Discrepancy} est conceptuellement cohérent avec la logique de surveillance
monétaire : une variation mensuelle anormale (ex. : effondrement soudain
des crédits à l'administration centrale ou une rupture dans les dépôts
transférables) est précisément caractérisée par l'impossibilité de
construire des associations non-triviales avec le reste de la série.

La longueur des séries (195 périodes) est suffisante pour pré-entraîner
un Anomaly Transformer de taille modeste. Les mécanismes d'attention
multi-têtes permettront de capturer simultanément les interactions
entre les 55 séries et les dépendances intra-série sur des horizons
variés (mensuel, trimestriel, annuel).

\textbf{Usage recommandé} : détection en temps quasi-réel des ruptures
dans les agrégats de premier niveau (Numéraire, Dépôts, Crédits,
Total des actifs), avec fenêtre glissante de 24 mois.

\subsection{Cohérence avec MKO-KAN}

\textbf{Applicabilité : Très élevée, avec adaptation au mode non supervisé.}

Le MKO-KAN présente une cohérence structurelle forte avec le dataset BEAC pour
plusieurs raisons complémentaires.

\textbf{Traitement séquentiel natif.}
La cellule TKANCell traite les séries temporelles de manière récurrente, ce
qui est naturellement aligné avec la fréquence mensuelle du dataset (195
périodes). Une fenêtre glissante de 12 à 24 mois peut être utilisée comme
entrée séquentielle, permettant de modéliser les dynamiques annuelles et les
effets de saisonnalité budgétaire.

\textbf{Interprétabilité réglementaire  atout décisif.}
Dans le contexte de la BEAC, la capacité du MKO-KAN à produire des formules
analytiques via la régression symbolique constitue un avantage différenciant
majeur par rapport aux approches boîte noire (AnoGAN, Anomaly Transformer).
Pour un auditeur ou un superviseur prudentiel, une règle de décision de la
forme $f(\text{Dépôts\_BQ} ,\, \text{Crédits\_ENF}) > \tau$ est directement
interprétable et opposable réglementairement.

\textbf{Modélisation des interactions entre indicateurs.}
Les nœuds MultKAN permettent d'encoder des interactions ciblées entre postes
du bilan  par exemple, la co-variation simultanée des crédits aux
sociétés non financières et des dépôts transférables non résidents,
caractéristique de certains chocs de financement externe.

\textbf{Détection de dérive structurelle.}
Le module de détection de dérive intégré au MKO-KAN (divergence de
Jensen-Shannon sur les distributions empiriques) peut signaler des ruptures
structurelles dans les agrégats, comme celles observées autour de 2020--2021
(choc COVID-19) et en 2024 (forte croissance des actifs non financiers).

\textbf{Condition d'adaptation.}
Le MKO-KAN est conçu pour la classification supervisée fraude/non-fraude.
Son application à la détection \textit{non supervisée} d'anomalies dans les
séries BEAC nécessite une adaptation : l'erreur de reconstruction de la
cellule T-KAN ou la déviation des poids d'arête par rapport à un modèle
entraîné sur la période de référence (2001--2019) peuvent servir de score
d'anomalie non supervisé.

\textbf{Usage recommandé} : détection d'anomalies sur les postes à haute
vélocité (crédits, dépôts transférables), avec production mensuelle de
rapports d'audit symboliques conformes aux exigences COBAC.

\subsection{Tableau récapitulatif}

\begin{longtable}{p{3.5cm} p{2cm} p{4cm} p{4cm}}
\toprule
\textbf{Méthode} & \textbf{Apt.} & \textbf{Atouts} & \textbf{Conditions} \\
\midrule
\endhead
LOF & Modérée & Non-paramétrique, interprétable & Réduction dim. préalable\\
Isolation Forest & Élevée & Scalable, robuste, non supervisé & Stationnarisation des séries\\
Autoencoder & Élevée & Modélisation temporelle & Peu de données (régularisation)\\
VAE & Élevée & Incertitude quantifiée & Intégrer attention (\textit{long-range})\\
AnoGAN / f-AnoGAN & Modérée & Distribution apprise & Taille dataset contraignante\\
Anomaly Transformer & Très élevée & Dépend. long terme, multiv. & Normalisation, pré-entraînement\\
MKO-KAN & Très élevée & Interprétable, temporel, interactions & Adaptation non supervisée requise\\
\bottomrule
\end{longtable}


% ═════════════════════════════════════════════════════════════════════════════
\section{Discussion et Recommandations}
\label{sec:discussion}
% ═════════════════════════════════════════════════════════════════════════════

L'analyse croisée entre les caractéristiques du dataset BEAC-Cameroun et
les méthodes IA disponibles suggère une stratégie en deux niveaux :

\begin{enumerate}
    \item \textbf{Niveau exploratoire} : l'Isolation Forest appliqué sur
          les résidus après décomposition (tendance + saisonnalité) fournit
          une baseline robuste, non supervisée et interprétable, adaptée à
          l'absence de labels d'anomalies certifiées.

    \item \textbf{Niveau avancé} : l'Anomaly Transformer ou un VAE avec
          attention permettront de capturer la complexité des interactions
          multivariées et des ruptures structurelles, notamment autour des
          épisodes de 2020--2021 et 2024. Une approche hybride combinant
          un Transformer pour la détection et un VAE pour la quantification
          de l'incertitude constitue la voie la plus prometteuse.
\end{enumerate}

Dans tous les cas, les étapes de prétraitement suivantes sont indispensables :
imputation des valeurs manquantes (interpolation linéaire ou méthode MICE),
normalisation par série (standardisation z-score), et décomposition STL
pour séparer tendance, saisonnalité et résidu.


% ═════════════════════════════════════════════════════════════════════════════
\section{Conclusion}
\label{sec:conclusion}
% ═════════════════════════════════════════════════════════════════════════════

Cette étude a présenté un état de l'art structuré des méthodes d'intelligence
artificielle pour la détection de variations atypiques, depuis le LOF
jusqu'à l'Anomaly Transformer, en passant par les autoencodeurs, les VAE
et les GAN. Chaque approche a été analysée à la lumière de ses fondements
théoriques et de ses caractéristiques opérationnelles.

L'application de ces méthodes aux données Monétaires issues
de la BEAC révèle une cohérence particulièrement forte avec l'Isolation
Forest, l'Anomaly Transformer et \textbf{MKO-KAN}. Ce dernier, développé
dans le cadre de ce stage, apporte une dimension supplémentaire : la
capacité à produire des formules analytiques auditables, alignée avec
les obligations de traçabilité algorithmique COBAC/BEAC. Le jeu de données,
riche de 55 séries temporelles mensuelles sur plus de deux décennies,
offre un terrain d'expérimentation prometteur pour la surveillance
macroprudentielle dans la zone CEMAC.


% ═════════════════════════════════════════════════════════════════════════════
% BIBLIOGRAPHIE
% ═════════════════════════════════════════════════════════════════════════════
\bibliographystyle{plain}

\begin{thebibliography}{99}

\bibitem{pang2021deep}
Pang, G., Shen, C., Cao, L., \& Van Den Hengel, A. (2021).
\textit{Deep Learning for Anomaly Detection: A Review}.
ACM Computing Surveys (CSUR), 54(2), 1--38.

\bibitem{breunig2000lof}
Breunig, M.M., Kriegel, H.P., Ng, R.T., \& Sander, J. (2000).
\textit{LOF: Identifying Density-Based Local Outliers}.
Proceedings of the 2000 ACM SIGMOD International Conference on
Management of Data, pp.~93--104.

\bibitem{liu2012isolation}
Liu, F.T., Ting, K.M., \& Zhou, Z.H. (2012).
\textit{Isolation-Based Anomaly Detection}.
ACM Transactions on Knowledge Discovery from Data (TKDD), 6(1), Article~3.
(Version initiale : ICDM 2008, pp.~413--422.)

\bibitem{hinton2006reducing}
Hinton, G.E., \& Salakhutdinov, R.R. (2006).
\textit{Reducing the Dimensionality of Data with Neural Networks}.
Science, 313(5786), 504--507.

\bibitem{zhou2017anomaly}
Zhou, C., \& Paffenroth, R.C. (2017).
\textit{Anomaly Detection with Robust Deep Autoencoders}.
Proceedings of the 23rd ACM SIGKDD International Conference on
Knowledge Discovery and Data Mining, pp.~665--674.

\bibitem{kingma2013auto}
Kingma, D.P., \& Welling, M. (2013).
\textit{Auto-Encoding Variational Bayes}.
arXiv preprint arXiv:1312.6114.

\bibitem{an2015vae}
An, J., \& Cho, S. (2015).
\textit{Variational Autoencoder Based Anomaly Detection Using
Reconstruction Probability}.
SNU Data Mining Center Technical Report.

\bibitem{schlegl2017unsupervised}
Schlegl, T., Seeböck, P., Waldstein, S.M., Schmidt-Erfurth, U., \& Langs, G. (2017).
\textit{Unsupervised Anomaly Detection with Generative Adversarial Networks
to Guide Marker Discovery}.
International Conference on Information Processing in Medical Imaging (IPMI),
pp.~146--157.

\bibitem{schlegl2019fanogan}
Schlegl, T., Seeböck, P., Waldstein, S.M., Langs, G., \& Schmidt-Erfurth, U. (2019).
\textit{f-AnoGAN: Fast Unsupervised Anomaly Detection with Generative
Adversarial Networks}.
Medical Image Analysis, 54, 30--44.

\bibitem{xu2022anomaly}
Xu, J., Wu, H., Wang, J., \& Long, M. (2022).
\textit{Anomaly Transformer: Time Series Anomaly Detection with Association
Discrepancy}.
International Conference on Learning Representations (ICLR 2022, Spotlight).

\bibitem{tanogan2020}
Bashar, M.A., \& Nayak, R. (2020).
\textit{TAnoGAN: Time Series Anomaly Detection with Generative Adversarial
Networks}.
arXiv preprint arXiv:2008.09567.

\bibitem{shap2025}
(Collectif). (2025).
\textit{Inter-layer Explainable Variational Autoencoder Model for Multivariate
Time Series Anomaly Detection}.
ScienceDirect, ScienceDirect, doi:10.1016/j.eswa.2025.XXX.

\bibitem{mamba2025}
(Collectif). (2025).
\textit{Mamba Adaptive Anomaly Transformer with Association Discrepancy
for Time Series}.
ScienceDirect, doi:10.1016/j.engappai.2025.XXX.

\bibitem{liu2024kan}
Liu, Z., Wang, Y., Vaidya, S., Ruehle, F., Halverson, J., Soljačić, M.,
Hou, T.Y., \& Tegmark, M. (2024).
\textit{KAN: Kolmogorov-Arnold Networks}.
arXiv preprint arXiv:2404.19756.

\bibitem{yanglu2024kan}
Lu, Y., \& Zhan, F. (2024).
\textit{Kolmogorov-Arnold Networks in Fraud Detection: Bridging the Gap
between Theory and Practice}.
Preprint. University of Cincinnati \& Stanford University.
5 septembre 2024.

\bibitem{li2024fastkan}
Li, Z. (2024).
\textit{FastKAN: Very Fast Implementation of Kolmogorov-Arnold Networks
(KAN)}.
GitHub repository. \texttt{https://github.com/ZiyaoLi/fast-kan}

\bibitem{zhou2024kanad}
Zhou, J. \textit{et al.} (2024).
\textit{KAN-AD: Anomaly Detection with Kolmogorov-Arnold Networks and
Fourier Series Edges}.
Preprint.

\bibitem{haosol2024meta}
Hao, J.-K., \& Solnon, C. (2024).
\textit{Méta-heuristiques et intelligence artificielle}.
Chapitre de livre. Université d'Angers \& INSA Lyon.

\bibitem{moukoko2025mkan}
Moukoko Ekambi, G.M. (2025).
\textit{MKO-KAN: Mobile Money KAN Fraud Detector 
Détection de fraude Mobile Money par réseaux de Kolmogorov-Arnold hybrides
en zone CEMAC}.
Rapport de stage, Banque des États de l'Afrique Centrale (BEAC).
ENSPD / Université de Douala, LIDIA Laboratory.

\end{thebibliography}

\end{document}